\documentclass[a4paper,12pt]{article}
\usepackage[T2A]{fontenc}
\usepackage[utf8]{inputenc}
\usepackage[russian]{babel}
\usepackage{amsmath,amssymb,mathtools}
\usepackage{helvet}
\renewcommand{\familydefault}{\sfdefault}
\usepackage{icomma}
\usepackage[a4paper,margin=23mm,top=23mm,bottom=22mm]{geometry}
\usepackage{microtype}
\usepackage{tikz}
\usetikzlibrary{calc,arrows.meta,positioning,shapes.geometric}
\usepackage{tabularx,booktabs,array}
\usepackage{enumitem}
\usepackage{hyperref}
\usepackage{parskip}
\usepackage{fancyhdr}
\usepackage{setspace}
\usepackage{xcolor}
\definecolor{accent}{HTML}{008080}
\definecolor{darkaccent}{HTML}{004D4D}
\definecolor{textgray}{HTML}{333333}
\definecolor{mutedgray}{HTML}{6C7476}
\definecolor{rulegray}{HTML}{C8D3D2}
\hypersetup{colorlinks=true,linkcolor=darkaccent,urlcolor=darkaccent,pdfauthor={Лёвин Артём Александрович},pdftitle={Чек-лист — Экономические задачи ЕГЭ}}
\onehalfspacing
\setlength{\parindent}{0pt}
\setlist[itemize]{leftmargin=6mm,topsep=2pt,itemsep=2pt,parsep=0pt}
\setlist[enumerate]{leftmargin=7mm,topsep=3pt,itemsep=9pt,parsep=1pt}
\newcommand{\key}[1]{\textcolor{darkaccent}{\textbf{#1}}}
\newcommand{\microhead}[1]{\vspace{2mm}{\large\textcolor{darkaccent}{\bfseries #1}}\par\vspace{1mm}}
\newcommand{\topic}[1]{\vspace{4mm}{\Large\textcolor{darkaccent}{\bfseries #1}}\par\vspace{1mm}\textcolor{rulegray}{\rule{\linewidth}{0.35pt}}\vspace{2mm}}
\newcommand{\checkitem}[1]{$\square$\hspace{2mm}#1\par\vspace{1mm}}
\pagestyle{fancy}
\fancyhf{}
\fancyhead[L]{\footnotesize\textcolor{mutedgray}{Экономические задачи ЕГЭ}}
\fancyhead[R]{\footnotesize\textcolor{mutedgray}{Ксения Клыкова}}
\fancyfoot[C]{\footnotesize\textcolor{mutedgray}{\thepage}}
\renewcommand{\headrulewidth}{0.25pt}
\renewcommand{\footrulewidth}{0pt}
\setlength{\headheight}{15pt}
\begin{document}
\begin{titlepage}
\thispagestyle{empty}
\centering
\vspace*{11mm}
{\footnotesize\bfseries\textcolor{mutedgray}{М А Т Е М А Т И К А}}\par
\vspace{34mm}
{\fontsize{30}{35}\selectfont\bfseries\textcolor{darkaccent}{Чек-лист}}\par
\vspace{11mm}
{\fontsize{20}{26}\selectfont\bfseries Экономические задачи ЕГЭ:\par
 вклады и кредиты}\par
\vspace{11mm}
{\large Подготовка к ЕГЭ по профильной математике}\par
\vfill
{\Large\bfseries Ксения Клыкова}\par
\vspace{10mm}
{\large 6 ноября 2026 г.}\par
\vspace{24mm}
{\normalsize Лёвин Артём Александрович}\par
{\small эксклюзивно для Ксении Клыковой}\par
\vspace{8mm}
\begin{tikzpicture}
\draw[accent!60,line width=0.8pt] (0,0) .. controls (2.6,0.8) and (4.5,-0.8) .. (7.2,0)
  .. controls (9.5,0.7) and (11,-0.3) .. (12.6,0.18);
\end{tikzpicture}
\vspace*{5mm}
\end{titlepage}

\topic{1. Как распознать финансовую схему}
В экономических задачах сначала определите, \key{что происходит с деньгами} и \key{как устроены выплаты}. Слова из условия подсказывают формулу.
\renewcommand{\arraystretch}{1.4}
\begin{center}
\small
\begin{tabularx}{\linewidth}{@{}p{0.25\linewidth}X@{}}
\toprule
\textbf{Схема} & \textbf{Указание в условии}\\
\midrule
Вклад & Деньги размещают в банке, на них начисляют проценты; возможны пополнения.\\
Дифференцированный кредит & Остаток долга каждый год уменьшается \emph{на одну и ту же сумму}.\\
Аннуитетный кредит & Кредит погашают \emph{равными ежегодными платежами}.\\
\bottomrule
\end{tabularx}
\end{center}
Обозначения во всём пособии: $S_0$~--- первоначальная сумма, $p$~--- годовая ставка \emph{в долях единицы}, $k=1+p$~--- повышающий коэффициент, $n$~--- число лет. Например, при $15\%$ имеем $p=0{,}15$, $k=1{,}15$.

\topic{2. Вклад с пополнением}
Пусть вклад открыт на два года, а в конце первого года \emph{после начисления процентов} внесено ещё $A$. Тогда в конце второго года:
\[
\boxed{S_2=S_0k^2+Ak.}
\]
\textbf{Почему так?} Первоначальная сумма приносит проценты дважды, а пополнение~--- один раз:
\[
S_0\xrightarrow{\times k}S_0k\xrightarrow{+A}S_0k+A
\xrightarrow{\times k}(S_0k+A)k.
\]

\microhead{Пример. Найти процентную ставку}
В начале положили $750$~тыс. руб., в конце первого года добавили $450$~тыс. руб. Через два года на счёте оказалось $1487{,}7$~тыс. руб. Найдите годовую ставку.
\[
750k^2+450k=1487{,}7.
\]
Умножим уравнение на $10$:
\[
7500k^2+4500k-14877=0.
\]
Дискриминант $D=466560000$, $\sqrt D=21600$. Тогда
\[
k=\frac{-4500+21600}{15000}=1{,}14,
\qquad p=k-1=0{,}14=14\%.
\]
Второй корень $k=-1{,}74$ не подходит: повышающий коэффициент положителен. \key{Ответ: $14\%$.}

\topic{3. Как извлечь корень из большого числа}
Этот приём помогает быстро вычислить дискриминант без громоздких вычислений.
\begin{enumerate}[label=\arabic*.]
\item Ограничьте число соседними квадратами десятков.
\item По последней цифре выделите возможные последние цифры корня.
\item \emph{Обязательно} возведите кандидата в квадрат и проверьте равенство.
\end{enumerate}
Например, $2500<2916<3600$, поэтому $50<\sqrt{2916}<60$. Число оканчивается на $6$, значит, кандидаты~--- $54$ и $56$. Проверка:
\[
54^2=2916,\qquad 56^2=3136,
\quad\text{поэтому}\quad\boxed{\sqrt{2916}=54}.
\]
Для большого дискриминанта полезно заранее выделить полные квадраты:
\[
\begin{aligned}
466560000&=16\cdot2916\cdot10000,\\
\sqrt{466560000}&=4\cdot54\cdot100=21600.
\end{aligned}
\]
\textbf{Важно.} Близость к меньшему квадрату позволяет выдвинуть предположение, но окончательный ответ подтверждается возведением в квадрат.

\topic{4. Дифференцированный кредит}
При таком кредите \key{основной долг} каждый год уменьшается на одну и ту же сумму $S_0/n$. Проценты начисляют на долг перед очередным платежом. Поэтому общие ежегодные платежи постепенно уменьшаются.

Остатки перед начислением процентов:
\[
S_0,\quad \frac{n-1}{n}S_0,\quad\ldots,\quad \frac1nS_0.
\]
Всего возвращается $S_0$ основного долга и проценты на эти остатки:
\[
\boxed{S_{\text{общ}}=S_0\left(1+p\frac{n+1}{2}\right).}
\]
\microhead{Пример. Общая сумма выплат}
Кредит $2{,}5$ млн руб. выдан на $4$ года под $15\%$ годовых. Остаток долга каждый год уменьшается на одинаковую величину.
\[
\begin{aligned}
S_{\text{общ}}&=2{,}5\left(1+0{,}15\cdot\frac{4+1}{2}\right)\\
&=2{,}5\cdot1{,}375=3{,}4375\ \text{млн руб.}
\end{aligned}
\]
\key{Ответ: $3\,437\,500$~руб.}

\microhead{Проверка по годам}
Основной долг уменьшается ежегодно на $2{,}5/4=0{,}625$ млн руб.; проценты по годам равны $0{,}375$, $0{,}28125$, $0{,}1875$, $0{,}09375$ млн руб. Их сумма~--- $0{,}9375$ млн руб. Добавляем $2{,}5$ млн руб. основного долга.

\topic{5. Аннуитетный кредит}
Здесь \key{каждый ежегодный платёж одинаков} и равен $x$. При этом проценты начисляют на ещё не погашенный долг, поэтому доля погашаемого основного долга меняется.

После каждого года остаток долга пересчитывается по правилу:
\[
\boxed{S_{i}=kS_{i-1}-x.}
\]
Если долг полностью погашен за $3$ года, то
\[
0=S_0k^3-xk^2-xk-x,
\]
откуда получаем
\[
\boxed{x=\frac{S_0k^3}{k^2+k+1}},\qquad
\boxed{S_{\text{общ}}=3x}.
\]
Для двух равных ежегодных платежей аналогично:
\[
\boxed{x=\frac{S_0k^2}{k+1}},\qquad
\boxed{S_{\text{общ}}=2x}.
\]
\microhead{Пример. Три равных платежа}
Взяли $496{,}5$~тыс. руб. на три года под $10\%$ годовых и возвращают долг тремя равными ежегодными платежами.
\[
\begin{aligned}
k&=1{,}1,\\
x&=\frac{496{,}5\cdot1{,}1^3}{1{,}1^2+1{,}1+1}
=\frac{660{,}8415}{3{,}31}=199{,}65\ \text{тыс. руб.}
\end{aligned}
\]
\[
S_{\text{общ}}=3\cdot199{,}65=598{,}95\ \text{тыс. руб.}
=\boxed{598950\ \text{руб.}}
\]
\textbf{Проверка долга:} $496{,}5\cdot1{,}1-199{,}65=346{,}5$; далее $346{,}5\cdot1{,}1-199{,}65=181{,}5$; наконец, $181{,}5\cdot1{,}1-199{,}65=0$.

\topic{6. Вычисления и единицы измерения}
\begin{itemize}
\item \textbf{Проценты:} $10\%=0{,}1$, $15\%=0{,}15$. В формулу подставляют долю единицы.
\item \textbf{Скобки:} $S_0(k^2+k+1)$ и $S_0k^2+S_0k+S_0$ равны; $S_0k^3$ требует возведения $k$ в третью степень.
\item \textbf{Единицы:} $1$ млн руб. $=1\,000\,000$ руб.; $1$ тыс. руб. $=1000$ руб.
\item \textbf{Порядок счёта:} сначала вычислите небольшие множители и степени, затем числитель, знаменатель и частное.
\end{itemize}

\clearpage
\thispagestyle{fancy}
\begin{center}
{\Large\bfseries\textcolor{darkaccent}{Алгоритм: как выбрать финансовую схему}}
\end{center}
\vspace{11mm}
\begin{center}
\begin{tikzpicture}[
>=Stealth,
node distance=12mm and 16mm,
base/.style={draw=accent!80,thick,align=center,font=\small,text width=44mm,inner sep=4.5mm},
act/.style={base,rounded corners=2mm},
ques/.style={base,diamond,aspect=2.5,text width=40mm,inner sep=1.5mm},
term/.style={base,ellipse,text width=40mm,inner sep=3mm},
arr/.style={->,thick,accent!85},
lab/.style={font=\footnotesize,fill=white,inner sep=1pt}
]
\node[term] (start) {Прочитайте условие\\и выделите\\движение денег};
\node[ques,below=12mm of start] (q1) {Сумму размещают\\на счёте и начисляют\\проценты?};
\node[act,right=6mm of q1] (a1) {\textbf{Вклад}\\Учитывайте сроки\\начисления процентов\\и пополнений};
\node[ques,below=14mm of q1] (q2) {Остаток долга\\уменьшается на равную\\сумму ежегодно?};
\node[act,right=6mm of q2] (a2) {\textbf{Дифференцированный}\\\textbf{кредит}\\Суммируйте долг\\и проценты на остатки};
\node[ques,below=14mm of q2] (q3) {По условию платежи\\равны между собой?};
\node[act,right=6mm of q3] (a3) {\textbf{Аннуитетный кредит}\\Найдите ежегодный\\платёж $x$, затем $nx$};
\node[act,below=12mm of q3] (other) {Уточните график:\\какая именно величина\\остаётся постоянной?};
\draw[arr] (start)--(q1);
\draw[arr] (q1.east)--node[lab,above]{да}(a1.west);
\draw[arr] (q1.south)--node[lab,right]{нет}(q2.north);
\draw[arr] (q2.east)--node[lab,above]{да}(a2.west);
\draw[arr] (q2.south)--node[lab,right]{нет}(q3.north);
\draw[arr] (q3.east)--node[lab,above]{да}(a3.west);
\draw[arr] (q3.south)--node[lab,right]{нет}(other.north);
\end{tikzpicture}
\end{center}
\clearpage

\begingroup
\setstretch{1.12}
\setlist[enumerate]{itemsep=4pt,topsep=2pt,parsep=0pt}
\topic{7. Тренировочный блок: Путь А}
\textcolor{mutedgray}{Выполните 10 заданий по возрастанию сложности. В задачах с процентами начисление происходит раз в год; платежи и пополнения вносятся в указанный момент. Ответы приведены на отдельной странице.}

\microhead{Базовый уровень}
\begin{enumerate}[label=\textbf{\arabic*.}]
\item Определите схему в каждом случае: \textbf{а)} банк начисляет проценты на вклад, в конце года его пополняют; \textbf{б)} остаток долга ежегодно сокращается на одну и ту же сумму; \textbf{в)} долг возвращают тремя равными ежегодными платежами.
\item На вклад положили $200$ тыс. руб. под $10\%$ годовых. В конце первого года после начисления процентов добавили $100$ тыс. руб. Сколько рублей будет на счёте в конце второго года?
\item На вклад положили $400$ тыс. руб., в конце первого года после начисления процентов добавили $120$ тыс. руб. Через два года получили $667$ тыс. руб. Найдите годовую процентную ставку.
\item Вклад $600$ тыс. руб. открыт на два года под $10\%$ годовых. В конце первого года после начисления процентов внесли ещё $A$ тыс. руб. В конце второго года на счёте оказалось $946$ тыс. руб. Найдите $A$.
\end{enumerate}
\microhead{Средний уровень}
\begin{enumerate}[label=\textbf{\arabic*.},start=5]
\item Найдите $\sqrt{7056}$, используя ближайшие квадраты десятков и последнюю цифру числа. Проверьте результат возведением в квадрат.
\item Кредит $1{,}5$ млн руб. взят на $4$ года под $15\%$ годовых. Каждый год остаток основного долга уменьшается на одну и ту же сумму. Найдите общую сумму выплат в рублях.
\item Кредит $1{,}8$ млн руб. погашают за $5$ лет по дифференцированной схеме. Общая сумма выплат равна $2{,}61$ млн руб. Найдите годовую процентную ставку.
\end{enumerate}
\microhead{Повышенный уровень}
\begin{enumerate}[label=\textbf{\arabic*.},start=8]
\item Кредит $210$ тыс. руб. взят на два года под $10\%$ годовых. Долг возвращают двумя равными ежегодными платежами. Найдите общую сумму выплат в рублях.
\item Кредит $331$ тыс. руб. взят на три года под $10\%$ годовых. Долг погашают тремя равными ежегодными платежами. Найдите общую сумму выплат в рублях.
\item Сравните общую сумму выплат по двум кредитам на $420$ тыс. руб. сроком на два года под $10\%$ годовых: \textbf{а)} остаток долга ежегодно сокращается на одинаковую сумму; \textbf{б)} предусмотрены два равных ежегодных платежа. На сколько рублей отличаются суммы выплат?
\end{enumerate}
\endgroup
\clearpage

\topic{8. Ответы}
\vspace{2mm}
\renewcommand{\arraystretch}{1.55}
\begin{center}
\begin{tabularx}{\linewidth}{@{}p{16mm}X@{}}
\toprule
\textbf{№} & \textbf{Ответ}\\
\midrule
1 & а) вклад; б) дифференцированный кредит; в) аннуитетный кредит.\\
2 & $352\,000$ руб.\\
3 & $15\%$.\\
4 & $A=200$ тыс. руб.\\
5 & $84$.\\
6 & $2\,062\,500$ руб.\\
7 & $15\%$.\\
8 & $242\,000$ руб.\\
9 & $399\,300$ руб.\\
10 & Дифференцированный: $483\,000$ руб.; аннуитетный: $484\,000$ руб. Разница $1000$ руб.\\
\bottomrule
\end{tabularx}
\end{center}
\vspace{8mm}
\microhead{Итоговая памятка}
\checkitem{По формулировке условия определяю вклад, дифференцированный или аннуитетный кредит.}
\checkitem{Перевожу процентную ставку в доли единицы: $p=r/100$, $k=1+p$.}
\checkitem{Для вклада учитываю, сколько лет работает каждый взнос.}
\checkitem{Для кредита различаю постоянное сокращение \emph{долга} и постоянный размер \emph{платежа}.}
\checkitem{Большие корни проверяю возведением в квадрат.}
\checkitem{Перед записью ответа перевожу тысячи или миллионы в рубли.}
\vfill
\begin{center}
\textcolor{mutedgray}{\footnotesize Лёвин Артём Александрович эксклюзивно для Ксении Клыковой}
\end{center}
\end{document}
